\section{总结与延伸}

Week 6 聚焦 AI 时代的软件测试与安全，呈现了一幅``双刃剑''图景：AI 既提升了安全测试的效率和覆盖面，也引入了全新的攻击向量。

\subsection{关键要点}

\begin{itemize}
    \item SAST（白盒）、DAST（黑盒）、SCA（依赖分析）三种技术互补
    \item AI Agent 引入五种新型攻击向量：prompt injection、tool misuse、intent breaking、identity spoofing、code attacks
    \item AI SAST 假阳性率 50--100\%，是当前最大的实用性障碍
    \item 非确定性输出使``完整覆盖''成为开放问题
    \item ``Shift left'' 安全理念在 AI 时代更易实践
    \item AI 生成代码的安全责任归属是未回答的伦理问题
\end{itemize}

\subsection{团队落地建议}

\begin{enumerate}
    \item 先把传统安全基线做好，再引入 AI 辅助检查，避免把 AI 当成“缺失安全工程”的替代品
    \item 为 AI 安全告警设置分级和升级路径，防止高假阳性直接压垮团队注意力
    \item 对 AI 生成的安全补丁保留更严格的 review 和测试门槛，特别是认证、权限和支付链路
\end{enumerate}

\subsection{事故响应视角下的 AI 安全}

Week 6 还有一个很实用的启发：安全不只是预防，更是响应。对于已经接入 Agent 的团队，事故处理流程也需要调整：

\begin{itemize}
    \item 先确认问题来自传统漏洞、提示注入还是工具滥用
    \item 保留 Agent 的执行轨迹、prompt、工具调用日志和环境状态
    \item 区分“模型误判”与“系统权限设计错误”，两者的修复路径完全不同
    \item 把事故复盘重新写回规则、测试和权限边界，而不是只修一个表面 bug
\end{itemize}

\begin{warningbox}{如果没有日志，很多 AI 安全事故几乎不可复盘}
传统 Web 服务至少有请求日志、错误堆栈和数据库审计；但 Agent 系统如果不记录 prompt、工具输入输出和关键中间决策，事后往往只能看到错误结果，却不知道它为什么做出这个动作。这会让修复停留在猜测层面。
\end{warningbox}

\subsection{安全策略总览表}

\begin{table}[H]
\centering
\begin{tabular}{p{0.18\textwidth} p{0.25\textwidth} p{0.31\textwidth}}
\toprule
风险面 & 主要防护手段 & 课程给出的原则 \\
\midrule
传统代码漏洞 & SAST / DAST / SCA 组合 & 不依赖单点检测，尽量前移到开发早期 \\
Agent 提示攻击 & Prompt 过滤、权限隔离、计划验证 & 把 prompt 当成新的攻击输入面来设计防护 \\
工具滥用与身份冒充 & 最小权限、身份校验、审计日志 & 不把工具调用默认视为可信动作 \\
AI 安全误报 & 分级告警、人工复核、回放验证 & 让 AI 提高召回率，但不独占最终裁决权 \\
\bottomrule
\end{tabular}
\caption{Week 6 的风险与防护对应关系}
\end{table}

\subsection{拓展阅读}

\begin{itemize}
    \item Semgrep 官网：\url{https://semgrep.dev}
    \item OWASP Top 10：\url{https://owasp.org/www-project-top-ten/}
    \item OWASP LLM Top 10：\url{https://owasp.org/www-project-top-10-for-large-language-model-applications/}
    \item Prompt Injection 攻防综述：Simon Willison's Blog
\end{itemize}

%% --- Fin del contenido --- %%

\end{document}
